\section{What a Provider Could Report}\label{sec:opportunity}

The set of results a provider could publish is not fixed, and it grows. When we count the benchmarks in our panel by the year each was published, the dated stock rises
from 11 at the end of 2019 to 62 by 2026. A benchmark that enters that stock stays
available to every release after it, so the denominator of any omission rate is
weakly increasing by construction.

What a single release could report grows more slowly and less evenly. Averaged by half-year, the number of eligible benchmarks per release runs 6.8 in
the first half of 2023 and falls to 3.9 in the first half of 2024, a period in which releases outpaced the instruments now used to score them. By the first half of 2026 it reaches 10.1. The two series move together over the sample and apart within it,
which is consistent with the opportunity set being shaped partly by release timing rather than by provider choice.

A count of omissions therefore measures the calendar. Unless reported tables
expand one for one with the stock, omission rates drift upward for arithmetic
reasons alone, and comparing early releases with late ones on that measure
compares eras rather than providers. The question that survives is composition:
which benchmarks are dropped, conditional on availability, against what an
indifferent selection over the same set would have produced. Answering it requires
knowing not only what was available but how the release stood on each available
benchmark, and it is that second requirement, not the first, that the rest of the
paper shows to be the hard one \citep{dranove2010}.
